\documentclass[titlepage]{article}

\usepackage{amsmath}
\usepackage{amsthm}
\usepackage{amssymb}
\usepackage{mdframed}
\usepackage{slantsc}
\usepackage[T1]{fontenc}
\usepackage{lmodern}
\usepackage{silence}
\WarningFilter{LaTeX}{Font shape}

\theoremstyle{remark}
\newmdtheoremenv{definition}{\textbf{Definition}}[section]
\newmdtheoremenv{proposition}[definition]{\textbf{Proposition}}
\newmdtheoremenv{lemma}[definition]{\textbf{Lemma}}
\newmdtheoremenv{theorem}[definition]{\textbf{Corollary}}
\newtheorem{example}[definition]{\textbf{Example}}

\renewcommand{\labelenumi}{(\roman*)}

\newcommand{\idealof}{\vartriangleleft}
\newcommand{\idealgen}[1]{\ensuremath{\left\langle#1\right\rangle}}
\newcommand{\seq}[2]{\ensuremath{#1_1,\dots,#1_{#2}}}

\title{Algebraic Geometry Notes}
\author{Alex Corbett}
\begin{document}
\date{\today}
\maketitle
\pagebreak
\section{Affine Algebraic Varieties}
  \begin{definition}
    Let $K$ be a field. Then an Affine Algrbraic Variety is a set of the form
    \[\mathcal{V}(I)=\{(a_1,a_2,\dots,a_n)\in K^n\mid f(a_1,a_2,\dots,a_n)=0, \forall f\in I\},\]
    where $I$ is an ideal $I\vartriangleleft K[x_1,x_2,\dots,x_n]$.
  \end{definition}
  \begin{proposition}
    Let $I_1,I_2$ be ideals of $K[x_1,\dots,x_n]$ be such that $I_1\subseteq I_2$, then $\mathcal{V}(I_1)\supseteq\mathcal{V}(I_2)$.
  \end{proposition}
  \begin{example}$ $
    \begin{enumerate}
      \item If $I=\{0\}$, then $\mathcal{V}(I)=K^n$.
      \item If $I=K[x_1,x_2,\dots,x_n]$, then $\mathcal{V}(I)=\{\emptyset\}$.
      \item If $I=\langle x^2+y^2-1\rangle\vartriangleleft\mathbb{R}[x,y]$, then $\mathcal{V}(I)$ is the circle of radius $1$ centred at the origin defined by the equation $x^2+y^2=1$ in $\mathbb{R}^2$.
      \item Let $f_1,f_2,\dots,f_r$ be polynomials of degree $1$ in $x_1,x_2,\dots,x_n$, and let $I=\langle f_1,f_2,\dots,f_r\rangle$ be the ideal generated byu them. Then $\mathcal{V}(I)$ is the set of solutions of the system of linear equations $f_1=f_2=\cdots=f_r=0$. Such a variety is called an affine subspace of $K^n$. This, in particular, includes sets containing a single point. If $P=(a_1,a_2,\dots,a_n)\in K^n$, then the ideal $I=\langle x_1-a_1,x_2-a_2,\dots,x_n-a_n\rangle$ defines the variety $\mathcal{V}(I)=\{P\}$.
      \item Let $I=\langle y^2-x^3\rangle\vartriangleleft\mathbb{R}[x,y]$, then $\mathcal{V}(I)$ is the cuspoidal cubic curve defined by $y^2=x^3$, also know as Neile's semicubal parabol, then $\mathcal{V}(I)$ is the cuspoidal cubic curve defined by $y^2=x^3$, also know as Neile's semicubal parabol, then $\mathcal{V}(I)$ is the cuspoidal cubic curve defined by $y^2=x^3$, also know as Neile's semicubal parabola.
      \item Let $I=\langle y^2-x^3-x^2\rangle\mathbb{R}[x,y]$, then $\mathcal{V}(I)$ is the nodal cubic curve defined by the equation $y^2=x^3+x^2$.
      \item If $I=\langle x,y\rangle\vartriangleleft\mathbb[x,y]$, then $\mathcal{V}(I)=\{(0,0)\}\subset\mathbb{R}^2$.
      \item If $I=\langle x^2,y^2\rangle\vartriangleleft\mathbb[x,y]$, then $\mathcal{V}(I)=\{(0,0)\}\subset\mathbb{R}^2$.
      \item If $I=\idealgen{x^2+y^2+z^2-1,x^2-x+y^2}\idealof\mathbb{R}[x,y,z]$, then $\mathcal{V}(I)\subset\mathbb{R}^3$ is a figure of eight curve from an intersection of the cylinder, $y^2=x^2-x$, and the sphere of radius $1$ at the origin.
      \item The special linear group $SL(n\mathbb{R})$ is also an affine algebraic variety. The set of all $n\times n$ real matrices is $\mathbb{R}^{n^2}$, the determinant is a polynomial in the entries of the matrix, and therefore $SL(n,\mathbb{R})=\mathcal{V}(\idealgen{\det{(A)}-1})$.
    \end{enumerate}
  \end{example}
  \begin{proposition}
    Let $K$ be a field. Let $f_1,f_2,\dots,f_r\in K[x_1,x_2,\dots,x_n]$ and let $I=\idealgen{f_1,f_2,\dots,f_r}$. Then
    \[\mathcal{V}(I)=\{(a_1,a_2,\dots,a_n)\in K^n\mid f_i(a_1,a_2,\dots,a_n)=0, \forall 1\le i\le r\}.\]
  \end{proposition}
  \begin{proof}
    We know $f_i\in I$ for each $i$, so if $(a_1,\dots,a_n)\in\mathcal{V}(I)$, then $f_i(a_1,\dots,a_n)=0$ for every $i$.
    Conversely, assume that $f_i(a_1,\dots,a_n)=0$ for every $i$. Any $f\in I$ can be written as a $K[x_1,\dots,x_n]$-linear combination of the $f_i$, so
    \[f(a_1,\dots,a_n)=\sum_{i=1}^rf_i(a_1,\dots,a_n)g_i(a_1,\dots,a_n)=0,\]
    so $(\seq{a}{n})$.
  \end{proof}
  \begin{definition}
    A ring $R$ is called Noetherian if all its ideals can be generated by finitely many elements.
  \end{definition}
  \begin{theorem}[Hilbert Basis Theorem]
    If $K$ is a field, then the polynomial ring $K[\seq{x}{n}]$ is Noetherian for any $n\ge0$.
  \end{theorem}
  \begin{example}
    \begin{enumerate}$ $
      \item If $R$ is Noetherian and $I\vartriangleleft R$, then $R/I$ is also Noetherian.
      \item $\mathbb{Z}$, the ring of integers, is Noetherian because it is a PID, so every ideal is generated by one element.
      \item $K[x_1,x_2,\dots]$, the polynomial rinf over the field $K$ in finitely many variables is an example of a non-Noetherian ring, because the ideal $\idealgen{x_1,x_2,\dots}$ is not finitely generated.
    \end{enumerate}
  \end{example}
  \begin{proposition}$ $
    \begin{enumerate}
      \item Let $V_1=\mathcal{V}(I_1),\dots,V_k=\mathcal{V}(I_k)$ be affine algebraic varieties in $K^n$. Then
        \[V_1\cup\cdots\cup V_k=\mathcal{V}(I_1\cap\cdots\cap I_k)=\mathcal{V}(I_1I_2\cdots I_k)\]
        is also an affine algebraic variety.
      \item Let $V_\alpha=\mathcal{V}(I_\alpha),\alpha\in A$ be affine algebraic varieties in $K^n$. Then
        \[\bigcap_{\alpha\in A}V_\alpha=\mathcal{V}\left(\sum_{\alpha\in A}I_\alpha\right)\]
        is also an affine algebraic variety.
      \item Let $V_1=\mathcal{V}(I_1)\subseteq K^m, V_2=\mathcal{V}(I_2)\subseteq K^n$ be affine algebraic varieties. Then $V_1\times V_2\subseteq K^m\times K^n\cong K^{m+n}$ is also an affine algebraic variety.
    \end{enumerate}
  \end{proposition}
  \begin{proof}$ $
    \begin{enumerate}
      \item 
        We first assume that $k=2$. Let $P\in V_1\cup V_2$ and let $f\in I_1\cap I_2$. If $P\in V_1$, then $f(P)=0$ because $f\in I_1$, and similarly if $P\in V_2$ then $f(P)=0$ because $f\in I_2$. Therefore in either case we have $f(P)=0$, so $P\in\mathcal{V}(1_1\cap I_2)$, so $V_1\cup V_2\subseteq\mathcal{V}(I_1\cap I_2)$.\\
        We also know $I_1I_2\subseteq I_1\cap I_2$, so we also have $\mathcal{V}(I_1\cap I_2)\subseteq\mathcal{V}(I_1I_2)$.\\
        Let now $P\in K^n\setminus(V_1\cup V_2)$. 
        Then $P\notin V_1$, so there exists $f_1\in I_1$ such that $f_1(P)\ne0$. Similarly, $P\ne V_2$ so there exists $f_2\in I_2$ such that $f_2(P)\ne0$. 
        Then $f_1f_2\in I_1I_2$ and $(f_1 f_2)(P)=f_1(P)f_2(P)\ne0$, therefore $P\notin\mathcal{V}(I_1I_2)$. 
        This implies $\mathcal{V}(I_1I_2)\subseteq V_1\cup V_2$. Hence we obtain
          \[V_1\cup V_2=\mathcal{V}(I_1\cap I_2)=\mathcal{V}(I_1I_2).\]
        For $k>2$ we use induction on $k$.

    \end{enumerate}
  \end{proof}
\end{document}
